\chapter{Gazebo Simulation and Gripper}
\label{ch:gazebo}

\section{World and physics}
\file{gp7_ros2/gp7_gazebo/worlds/pick_place.sdf} defines the
\code{gp7_pick_place} world. It loads physics, user commands, and scene
broadcasting systems; applies gravity $(0,0,-9.81)$ m/s$^2$; and configures a
nominal maximum physics step of \SI{0.001}{s} and real-time factor 1. The
selected physics engine was observed to warn that mimic constraints are not
supported. The environment includes a directional light and static ground
plane.

Box A is a static blue box of size $(0.30,0.30,0.20)$ m with its model center
at $(0.55,0.35,0.10)$ m. This corresponds to the ground-referenced position
$(0.55,0.35,0)$ m used in the documentation. Box B is a dynamic red cube of
side \SI{0.05}{m}, mass \SI{0.05}{kg}, and center height \SI{0.025}{m} when on
the ground. The world file starts it at $(0.55,-0.20,0.025)$ m; the task
launch creates a temporary edited world for either far or seeded-near pickup.
Gazebo's pose publisher emits model pose vectors at \SI{50}{Hz}.

\section{Robot spawn and control coupling}
The task launch calls the \code{ros_gz_sim create} executable with the
generated URDF string. The robot description contains the external
\code{gz_ros2_control/GazeboSimSystem} plugin and controller YAML path.
Gazebo then exposes simulated joint state/command interfaces to the ROS 2
controller manager. Controller actions and gripper command topics alter the
Gazebo robot. The clock bridge aligns ROS timers and task waits with simulated
time.

\section{Finger actuation and mimic limitation}
Only \code{gripper_driver} is commanded by the task. The URDF describes five
dependent finger revolute joints with mimic multipliers of $+1$ or $-1$. The
task's \code{fingers(target)} function interpolates from the observed current
driver angle to target using quintic scaling; its nominal duration is
\begin{equation}
 T_g=\max\left(1\si{s},\frac{|q_{target}-q_{start}|}{0.25\si{rad/s}}\right).
\end{equation}
It publishes at approximately \SI{20}{ms} simulation-time steps, waits an
additional \SI{0.4}{s}, and rejects if measured driver angle differs by more
than 0.02 rad. During task grasp the close target is $-0.162$ rad; the fully
open target is $+0.15$ rad. The manual wrapper allows commands down to
$-0.5$ rad, but the task does not use that full-close setting.

The physics engine warning means URDF mimic metadata did not establish
realistic finger coupling in the recorded simulation. A visually moving
driver or a successfully completing task is not evidence that both fingers
generated accurate, symmetric contact forces.

\section{Idealized object attachment}
The task-specific function \code{task_description} appends a Gazebo
\code{gz-sim-detachable-joint-system} plugin to the URDF. It identifies
\code{link_6_t} as parent, model \code{box_b} and link \code{link} as child,
and uses Gazebo Transport topics:
\begin{align*}
&\code{/gp7/grasp/attach},\qquad
\code{/gp7/grasp/detach},\\
&\code{/gp7/grasp/state}.
\end{align*}
The task node waits for state text \code{attached} or \code{detached}. Gazebo's
plugin is initially attached when a model starts, so the node explicitly asks
for and confirms detachment before approach. After the fingers close, the node
requests attachment. During transfer, it verifies that the attachment remains
confirmed and reads the cube pose. At placement it checks the object is close
to target, detaches, opens, retreats, and lets gravity settle the cube.

This fixed-joint mechanism forces the cube to follow the robot rigidly. It
does not calculate fingertip contact, friction cones, normal force, grasp
stability, tactile feedback, or slip. Consequently, reported micrometre-scale
simulation placement errors characterize the constrained simulation pose and
numerics; they are not force-based grasp accuracy or a physical robot claim.

\section{Collision and contact scope}
Gazebo has visual/collision objects and gravity, and the dynamic cube settles
after detach. However, the planner does not query collision geometry or plan
around arbitrary obstacles. Its far route includes a hand-selected inward
waypoint to avoid cutting through the base; this is not general collision
checking. The gripper collision tags reference visual DAE geometry, so their
contact representation is also not independently validated as a realistic
finger-contact model.

\section{Grasp improvement path (future work)}
A physical or contact-realistic extension would need supported mimic control
or independent finger joints, calibrated contact/collision geometry, suitable
friction coefficients, validated inertias, explicit gripper force/current
limits, and experiments measuring retention and slip. The detachable joint
could remain a software test fixture but should be disabled for those tests.
None of these future changes are claimed as implemented.
